\par\clearpage
\SourceLeaf{147}\GaussVTheoryEarlyFolio{143}
\[
\begin{aligned}
X&=-\left(\frac{\mathrm{d}P'}{\mathrm{d}u}+\frac{\mathrm{d}P''}{\mathrm{d}u}+\frac{\mathrm{d}P'''}{\mathrm{d}u}+\text{u. s. w.}\right)\\
Y&=-\frac1{\sin u}\left(\frac{\mathrm{d}P'}{\mathrm{d}\lambda}+\frac{\mathrm{d}P''}{\mathrm{d}\lambda}+\frac{\mathrm{d}P'''}{\mathrm{d}\lambda}+\text{u. s. w.}\right)\\
Z&=2P'+3P''+4P'''+\text{u. s. w.}
\end{aligned}
\]

\Article{20}
Verbinden wir nun mit diesen Sätzen das bekannte Theorem, dass jede Function von $\lambda$ und $u$, die für alle Werthe von $\lambda$ von $0$ bis $360^0$, und von $u$ von $0$ bis $180^0$ einen bestimmten endlichen Werth hat, in eine Reihe von der Gestalt
\[
P^0+P'+P''+P'''+\text{u. s. w.}
\]
entwickelt werden kann, deren allgemeines Glied $P^{\mathrm{n}}$ der obigen partiellen Differentialgleichung Genüge leistet, dass eine solche Entwicklung nur auf Eine bestimmte Art möglich ist, und dass diese Reihe immer convergirt, so erhalten wir folgende merkwürdige Sätze:

I. Die Kenntniss des Werths von $V$ in allen Punkten der Erdoberfläche reicht hin, um den allgemeinen Ausdruck von $V$ für den ganzen unendlichen Raum ausserhalb der Erdfläche daraus abzuleiten, und somit auch die Bestimmung der Kräfte $X,Y,Z$ nicht bloss auf der Erdoberfläche, sondern auch für den ganzen unendlichen Raum ausserhalb derselben. Offenbar ist dazu nur nöthig, $\frac VR$ nach dem erwähnten Theorem in eine Reihe zu entwickeln.

Es soll daher im Folgenden das Zeichen $V$ immer in der auf die Oberfläche der Erde beschränkten Bedeutung verstanden werden, wenn das Gegentheil nicht ausdrücklich gesagt ist, oder als diejenige Function von $\lambda$ und $u$, welche aus dem allgemeinen Ausdruck hervorgeht, wenn $r=R$ gesetzt wird, also
\[
V=R(P'+P''+P'''+\text{u. s. w.})
\]

II. Die Kenntniss des Werthes von $X$ in allen Punkten der Erdoberfläche reicht hin, um alles in I. angeführte zu erlangen. In der That ist nach Art. 15 das Integral $\int_0^u X\mathrm{d}u=\frac{V^0-V}{R}$, wenn $V^0$ den Werth von $V$ im Nordpole bedeutet, und die Entwicklung von $\int_0^u X\mathrm{d}u$ in eine Reihe der erwähnten Form muss nothwendig mit
\par\clearpage
\SourceLeaf{148}\GaussVTheoryEarlyFolio{144}
\[
V^0-P'-P''-P'''-\text{u. s. w.}
\]
identisch sein.

III. Auf gleiche Weise und unter Bezugnahme auf Art. 16 ist klar, dass die Kenntniss des Werthes von $Y$ auf der ganzen Erde verbunden mit der Kenntniss von $X$ in allen Punkten einer von einem Erdpole zum andern laufenden Linie zur Begründung der \textit{vollständigen} Theorie des Erdmagnetismus zureicht.

IV. Endlich ist klar, dass die vollständige Theorie auch aus der blossen Kenntniss des Werthes von $Z$ auf der ganzen Erdfläche abzuleiten ist. In der That, wenn $Z$ in eine Reihe entwickelt wird
\[
Z=Q^0+Q'+Q''+Q'''+\text{u. s. w.}
\]
so dass das allgemeine Glied der mehrerwähnten partiellen Differentialgleichung Genüge leistet, so wird nothwendig $Q^0=0$ und $P'=\frac12Q'$, $P''=\frac13Q''$, $P'''=\frac14Q'''$ u. s. w sein müssen.

\Article{21}
Wegen der einfachen Art der Abhängigkeit der einzelnen Kräfte $X,Y,Z$ von einer einzigen Function $V$, und des einfachen Zusammenhanges, in welchem jene unter sich stehen, sind dieselben weit mehr geeignet, zur Grundlage der Theorie zu dienen, als der gewöhnliche Ausdruck der magnetischen Kraft durch die drei Elemente ganze Intensität, Inclination und Declination, oder vielmehr, die letztere Art. so natürlich sie an sich scheint, wo es nur darauf ankommt die Thatsachen aufzufassen, kann unmittelbar gar nicht zur Begründung der Theorie, wenigstens nicht zur ersten Begründung, dienen, ehe sie nicht in die andere Form übersetzt ist. In dieser Beziehung wäre es daher sehr wünschenswerth, dass eine allgemeine graphische Darstellung der horizontalen Intensität veranstaltet würde, theils weil diese dem für die Theorie brauchbaren näher steht als die ganze Intensität, theils weil jene bei weitem in den meisten Fällen das ursprünglich wirklich beobachtete, die letztere hingegen nur durch Rechnung vermittelst der Inclination daraus abgeleitet ist. Die Elemente des horizontalen Magnetismus für sich rein zu erhalten, bleibt um so mehr zu empfehlen, da sie durch die gegenwärtigen Hülfsmittel sich mit überwiegender Schärfe bestimmen lassen, und man sollte wenigstens niemals mit Unterdrückung der beobachteten horizontalen
\par\clearpage
\SourceLeaf{149}\GaussVTheoryEarlyFolio{145}
\noindent Intensität die durch Rechnung daraus abgeleitete ganze Intensität bekannt machen, ohne die bei der Rechnung angewandte Inclination mit anzugeben, damit derjenige, welcher sie für die Theorie benutzen will, im Stande sei, die ursprünglichen Zahlen unverfälscht wieder herzustellen.

So interessant es übrigens auch sein würde, die ganze Theorie des Erdmagnetismus allein auf Beobachtungen der horizontalen Nadel zu gründen, und damit den verticalen Theil oder die Inclination zu anticipiren, so ist es doch dazu gegenwärtig noch viel zu früh: die Mangelhaftigkeit der jetzt zu Gebote stehenden Data verstattet nicht, auf den Mitgebrauch des verticalen Theils zu verzichten. Im Grunde empfängt auch die Theorie schon dadurch ihre Bestätigung, wenn die Vereinbarkeit sämmtlicher Elemente unter Ein Princip nachgewiesen werden kann.

\Article{22}
Wenn wir gleich \textit{a priori} gewiss sind, dass die Reihen für $V,X,Y,Z$ convergiren, so lässt sich doch im voraus nichts über den Grad der Convergenz bestimmen. Wären entweder die Sitze der magnetischen Kräfte auf einen mässigen Raum um den Mittelpunkt der Erde her beschränkt, oder fände eine solche Vertheilung der magnetischen Flüssigkeiten in der Erde Statt, die jenem Falle äquivalirte, so würden die Reihen sehr schnell convergiren müssen; je weiter hingegen jene Sitze bis gegen die Oberfläche hin sich erstrecken, und je unregelmässiger die Vertheilung ist, desto mehr wird man auf eine langsame Convergenz sich gefasst halten müssen.

Bei der praktischen Anwendung ist absolute Genauigkeit unerreichbar: man verlangt nur einen den Umständen angemessenen Grad von Annäherung. Je langsamer nun die Convergenz ist, eine desto grössere Anzahl von Gliedern wird berücksichtigt werden müssen, um einen bestimmten Grad von Genauigkeit zu erreichen.

Nun enthält $P'$ drei Glieder, und erfordert also die Kenntniss von drei Coëfficienten $g^{1,0},g^{1,1},h^{1,1}$; $P''$ erfordert fünf Coëfficienten, $P'''$ sieben, $P''''$ neun u. s. w. Indem wir also $P',P'',P'''$ u. s. w. als Grössen erster, zweiter, dritter Ordnung u. s. w. betrachten, erhellt, dass wenn die Rechnung bis zu den Grössen der Ordnung $n$ einschliesslich getrieben werden soll, die Werthe von $nn+2n$
\par\GaussVTheoryGathering{19}
\par\clearpage
\SourceLeaf{150}\GaussVTheoryEarlyFolio{146}
\noindent Coëfficienten ausgemittelt werden müssen, also z. B. 24, wenn man bis zur vierten Ordnung gehen will.

\expandafter\def\csname HistoricalNoteText18\endcsname{Am wenigsten nachtheilig würden bei einer solchen Bestimmungsweise diese Umstände einwirken, wenn die acht Punkte ganz symmetrisch auf der Erdoberfläche vertheilt waren, d. i. wenn sie mit den Ecken eines in der Erdkugel eingeschriebenen Würfels zusammenfielen, oder doch einer solchen Lage sehr nahe kämen.}
Jeder gegebene Werth von $X,Y$ oder $Z$, für gegebene Werthe von $u$ und $\lambda$ verschafft uns eine Gleichung zwischen den Coëfficienten, mithin geben die vollständig bekannten Elemente des Erdmagnetismus von jedem Orte drei Gleichungen. Dürfte man also annehmen, dass nur die Glieder bis zur vierten Ordnung merklich bleiben, so würden zur Bestimmung aller nöthigen Coëfficienten die vollständigen Beobachtungen von acht verschiedenen Punkten, theoretisch betrachtet, zureichen: allein jene Voraussetzung ist schwerlich zulässig, und so würden die allen Beobachtungen anhängenden zufälligen Fehler verbunden mit der Vernachlässigung der Glieder der höhern Ordnung die Eliminationsresultate sehr entstellen können\HistoricalNote{18}. Den schädlichen Einfluss dieser Umstände zu vermindern, müsste man eine viel grössere Anzahl von Beobachtungsstücken, als unbekannte Grössen sind, von weit auseinander liegenden Punkten aus allen Theilen der Erde, zum Grunde legen, und die unbekannten Grössen nach der Methode der kleinsten Quadrate daraus ableiten. So einförmig indessen, da alle Gleichungen nur linearisch sind, die Ausführung eines solchen Geschäfts auch sein würde, so möchte doch die ausserordentliche aus der grossen Menge der unbekannten Grössen und Gleichungen entspringende Weitläuftigkeit auch den muthigsten Rechner abschrecken, die Arbeit in dieser Form jetzt schon zu unternehmen, zumal da das Einschleichen von unzuverlässigen Beobachtungsstücken oder von Rechnungsfehlern den Erfolg ganz verderben könnte.

{\renewcommand{\Article}[1]{\par\medskip\Needspace{4\baselineskip}\pdfbookmark[1]{Art. #1}{erdmagnetismus-art-#1}\centerline{\textbf{#1}}\nobreak\smallskip}\Article{23}}
Es gibt aber ein anderes Verfahren, welches, von einem Theile dieser Schwierigkeiten frei, sich vorzugsweise für den ersten anzustellenden Versuch zu eignen scheint, und welches wir hier entwickeln wollen, ohne die Bedenklichkeiten zu verschweigen, denen die Anwendung desselben bei jetziger Lage der Sachen noch unterliegt. Dies Verfahren setzt die Kenntniss aller drei Elemente in Punkten voraus, die auf einer hinlänglichen Anzahl von Parallelkreisen so gruppirt sind,
\par\clearpage
\SourceLeaf{151}\GaussVTheoryEarlyFolio{147}
\noindent dass jeder Parallelkreis dadurch in eine hinlängliche Anzahl gleicher Stücke getheilt wird.

Aus den in gewöhnlicher Form gegebenen Elementen hat man zuvörderst die numerischen Werthe von $X,Y$ und $Z$ abzuleiten.

Man bringt sodann, nach bekannter Methode, die Werthe von $X,Y$ und $Z$ auf jedem Parallelkreise in die Form
\[
\begin{aligned}
X&=k+k'\cos\lambda+K'\sin\lambda+k''\cos2\lambda+K''\sin2\lambda\\&\quad+k'''\cos3\lambda+K'''\sin3\lambda+\text{u. s. w.}\\
Y&=l+l'\cos\lambda+L'\sin\lambda+l''\cos2\lambda+L''\sin2\lambda\\&\quad+l'''\cos3\lambda+L'''\sin3\lambda+\text{u. s. w.}\\
Z&=m+m'\cos\lambda+M'\sin\lambda+m''\cos2\lambda+M''\sin2\lambda\\&\quad+m'''\cos3\lambda+M'''\sin3\lambda+\text{u. s. w.}
\end{aligned}
\]
Man erhält also für jeden der Coëfficienten $k,l,m,k'$ u. s. w. so viele Werthe, als Parallelkreise behandelt sind.

Der Theorie zufolge sollte auf jedem Parallelkreise $l=0$ werden; die aus der Rechnung hervorgehenden Werthe von $l$ geben also schon eine Art von Maassstab für den Grad von Unzuverlässigkeit, welcher die zum Grunde gelegten Zahlen noch unterliegen.

Aus den Gleichungen
\[
\begin{aligned}
k&=-g^{1,0}\frac{\mathrm{d}P^{1,0}}{\mathrm{d}u}-g^{2,0}\frac{\mathrm{d}P^{2,0}}{\mathrm{d}u}-g^{3,0}\frac{\mathrm{d}P^{3,0}}{\mathrm{d}u}-\text{u. s. w.}\\
m&=2g^{1,0}P^{1,0}+3g^{2,0}P^{2,0}+4g^{3,0}P^{3,0}+\text{u. s. w.}
\end{aligned}
\]
deren Gesammtanzahl doppelt so gross ist, als die Anzahl der Parallelkreise, wird man, nachdem in $\frac{\mathrm{d}P^{1,0}}{\mathrm{d}u}$ u. s. w. und in $P^{1,0}$ u. s. w. die entsprechenden Zahlwerthe von $u$ substituirt sind, von den Coëfficienten $g^{1,0},g^{2,0},g^{3,0}$ u. s. w. so viele, als berücksichtigt werden sollen, nach der Methode der kleinsten Quadrate bestimmen.

Eben so dienen die Gleichungen
\[
\begin{aligned}
-k'&=g^{1,1}\frac{\mathrm{d}P^{1,1}}{\mathrm{d}u}+g^{2,1}\frac{\mathrm{d}P^{2,1}}{\mathrm{d}u}+g^{3,1}\frac{\mathrm{d}P^{3,1}}{\mathrm{d}u}+\text{u. s. w.}\\
L'&=g^{1,1}\frac{P^{1,1}}{\sin u}+g^{2,1}\frac{P^{2,1}}{\sin u}+g^{3,1}\frac{P^{3,1}}{\sin u}+\text{u. s. w.}\\
m'&=2g^{1,1}P^{1,1}+3g^{2,1}P^{2,1}+4g^{3,1}P^{3,1}+\text{u. s. w.}
\end{aligned}
\]
deren Anzahl zusammen dreimal so gross ist, als die Anzahl der Parallelkreise, zur Bestimmung der Coëfficienten $g^{1,1},g^{2,1},g^{3,1}$ u. s. w.; so wie folgende
\par\GaussVTheoryGathering{19\textsuperscript{*}}
